% Part: proof-theory
% Chapter: natural-deduction
% Section: grafting
%READERNOTE{SOL6-A300: कलम-संयोजनापूर्वी बाहेरच्या सिद्धतेच्या मुक्ती-खुणा आत बसवलेल्या सिद्धतेच्या उघड्या गृहीतकांच्या खुणांपासून वेगळ्या ठेवायच्या; अन्यथा समान सूत्र असले तरी गृहीतक अनवधानाने मुक्त होऊ शकते.}

\documentclass[../../../include/open-logic-section]{subfiles}

\begin{document}

\olfileid{pt}{nat}{gra}

\olsection{\usetoken{P}{proof} यांचे कलम-संयोजन}

~\Log{N1c} आणि~\Log{N1i}
मधील !!^{proof}s म्हणजे
गृहीतके~$!B$ यांपासून
!!{formula}s~$!A$ यांच्या
!!{proof}s होत. आता $!B$
चीही !!a{proof} आहे असे समजा.
$!B$ ची !!{proof} आणि
$!B$ पासून $!A$ ची !!{proof}
एकत्र करून, $!B$ हे गृहीतक
न वापरणारी $!A$ ची
!!a{proof} कशी मिळवता येईल?

एक मार्ग म्हणजे
~\Intro{\lif} लावून $!A$
च्या !!{proof} मधील
$!B$ हे गृहीतक काढून टाकणे.
मग मिळालेली $!B \lif !A$
ची !!{proof} आणि $!B$
ची !!{proof} एकत्र
केल्यास पुढील सिद्धता मिळते:
\[
\AxiomC{$\Discharge{!B}{x}$}
\DeduceC{$!A$}
\DischargeRule{\Intro{\lif}}{x}
\UnaryInfC{$!B \lif !A$}
\AxiomC{}
\DeduceC{$!B$}
\RightLabel{\Elim{\lif}}
\BinaryInfC{$!A$}
\DisplayProof
\]
हे सरळ आहे, पण काहीसे
अनावश्यक वळण घेते:
$\lif$ आणून लगेच त्याचे
निराकरण का करावे?
$!B \lif !A$ मधून जाणारा
असा “वळसा” टाळता यायला हवा.
(असे वळसे नेहमी टाळता
येतात, हाच प्रत्यक्षात
सामान्यरूपीकरण प्रमेयाचा आशय आहे.)

~\Log{N1c} आणि~\Log{N1i}
मध्ये या प्रसंगी तो वळसा
सहज टाळता येतो: $!B^x$
या गृहीतकांच्या जागी $!B$
ची संपूर्ण !!{proof} बसवता येते.
इतर सिद्धतांच्या गृहीतकांवर
!!{proof}s चे असे
“कलम-संयोजन” उपयुक्त आहे;
म्हणून त्याची व्याख्या नोंदवू.

\begin{defn}
समजा $\delta_1$ ही उघडे
गृहीतक~$!B^x$ यापासून
$!A$ ची \Log{N1c}-!!{proof}
(किंवा \Log{N1i}-!!{proof}) आहे,
आणि $\delta_2$ ही $!B$
ची \Log{N1c}-!!{proof}
(\Log{N1i}-!!{proof}) आहे.
मग $\delta_1$ मधील
$!B^x$ च्या प्रत्येक
निरसन न झालेल्या उदाहरणाच्या
जागी $\delta_2$ बसवून
मिळणारा परिणाम म्हणजे
$\Subst{\delta_1}{\delta_2}{!B^x}$.
\end{defn}

~$\delta_1$ आणि $\delta_2$
मध्ये एकाच खूणेची, पण वेगवेगळी
गृहीतक-!!{formula}s नसतील,
आणि $\delta_2$ च्या कोणत्याही
उघड्या गृहीतकात $\delta_1$
चे आयगेन चर येत नसेल,
आणि $\delta_2$ च्या उघड्या गृहीतकाची कोणतीही खूण
$\delta_1$ मधील मुक्ती-नियमाची खूण नसेल,
तर परिणामी
$\Subst{\delta_1}{\delta_2}{!B^x}$
ही योग्य !!{proof} आहे.
तिची उघडी गृहीतके म्हणजे
$\delta_1$ मधील बसवलेल्या
उदाहरणांव्यतिरिक्त
उरलेली उघडी गृहीतके,
आणि $\delta_2$ मधील
उघडी गृहीतके.
!!{proof}s चे कलम-संयोजन
करताना या अटी सहसा पूर्ण होतात.
~$\delta_1$ आणि $\delta_2$
या मोठ्या !!{proof}
च्या उपसिद्धता असतील,
तर पहिली अट पूर्ण होते.
दुसरी अट पूर्ण नसेल,
तर $\delta_1$ ची आयगेन चरे
पुन्हा नाव देऊन ती पूर्ण
करता येते.
ही खुणांची नवी अटही आवश्यक आहे: दोन्हीकडे तेच सूत्र
असले तरी आत बसवलेले उघडे गृहीतक बाहेरच्या मुक्ती-नियमाने
अनवधानाने मुक्त होता कामा नये. कलम-संयोजनापूर्वी
$\delta_1$ मधील मुक्ती-खुणा आणि त्यांना बांधलेली गृहीतके
एकत्र नवी नावे देऊन, $\delta_2$ च्या खुणांपासून वेगळी
करता येतात; मूळ उघड्या गृहीतकांची नावे जपायची.

तशीच व्याख्या
\Log{N2c} आणि~\Log{N2i}
यांनाही लागू होते.

\begin{defn}
समजा $\delta_1$ ही
$x:!B, \Gamma_1 \Sequent !A$
ची \Log{N2c}-!!{proof}
(किंवा \Log{N2i}-!!{proof}) आहे;
आणि $\delta_2$ ही
$\Gamma_2 \Sequent !B$
ची \Log{N2c}-!!{proof}
(\Log{N2i}-!!{proof}) आहे.
मग $\delta_1$ मधील
अंतिम क्रमवर्तीपर्यंत मुक्त न झालेल्या
$x:!B$ शी संबंधित प्रत्येक $x:!B \Sequent !B$
स्वयंसिद्धकाच्या जागी $\delta_2$ बसवून,
त्याखालील प्रत्येक क्रमवर्तीच्या संदर्भातील
संबंधित $x:!B$ काढून त्याऐवजी $\Gamma_2$ जोडून
मिळणारा परिणाम म्हणजे
$\Subst{\delta_1}{\delta_2}{x:!B}$.
\end{defn}

\begin{prop}
जर $\delta_1 \Proves
x:!B, \Gamma_1 \Sequent
!A$ आणि $\delta_2 \Proves
\Gamma_2 \Sequent !B$,
तर $\Subst{\delta_1}{\delta_2}{x:!B}
\Proves \Gamma_1, \Gamma_2 \Sequent
!A$, मात्र दोन्ही सिद्धतांतील गृहीतक-खुणा सुसंगत असाव्यात,
$\Gamma_2$ मधील कोणतीही खूण $\delta_1$ मधील
मुक्ती-नियमाची खूण नसावी आणि $\delta_1$ चे
कोणतेही आयगेन चर $\Gamma_2$ मध्ये येता कामा नये.
\end{prop}

\begin{proof}
  $\pheight{\delta_1}$ वर
  विगमन करून.
\end{proof}

\end{document}
